% Intervalos y valores p recalculados con la SEMILLA como unidad de
% remuestreo por final-hardening/recompute_clustered_ci.py.
% Bootstrap pareado sobre semillas, B=10000, semilla de analisis 20260919;
% Wilcoxon de una cola en la direccion predeclarada sobre diferencias
% por semilla, y Holm dentro de la familia de la campana. Los
% contrastes deterministas por semilla no llevan intervalo ni valor p.
\newcommand{\SPTwoWorlds}{1560}
\newcommand{\SPTwoComparisonRows}{20280}
\newcommand{\SPTwoAblationRows}{9360}
\newcommand{\SPTwoScoreCoverage}{0.514}
\newcommand{\SPTwoScoreSuccess}{0.487}
\newcommand{\SPTwoScoreRuntime}{40.84}
\newcommand{\SPTwoCoverageCoverage}{0.529}
\newcommand{\SPTwoCoverageSuccess}{0.227}
\newcommand{\SPTwoCoverageRuntime}{11.31}
\newcommand{\SPTwoNeuralGap}{0.109}
\newcommand{\SPTwoNeuralSuccess}{0.263}
\newcommand{\SPTwoNeuralRuntime}{1.01}
\newcommand{\SPTwoNeuralParameters}{121}
\newcommand{\SPTwoSmithGap}{0.343}
\newcommand{\SPTwoSmithSuccess}{0.103}
\newcommand{\SPTwoSmithRuntime}{0.35}
\newcommand{\SPTwoPlainSuccess}{0.051}
\newcommand{\SPTwoMarginalSuccess}{0.135}
\newcommand{\SPTwoPlainGap}{0.372}
\newcommand{\SPTwoMarginalGap}{0.157}
\newcommand{\SPTwoPlainIncomplete}{0.946}
\newcommand{\SPTwoMarginalIncomplete}{0.848}
\newcommand{\SPTwoPlainAlignment}{0.049}
\newcommand{\SPTwoMarginalAlignment}{0.137}
\newcommand{\SPTwoGapEffect}{-0.2145}
\newcommand{\SPTwoGapCILow}{-0.2230}
\newcommand{\SPTwoGapCIHigh}{-0.2065}
\newcommand{\SPTwoGapPHolm}{4{,}55\times10^{-12}}
\newcommand{\SPTwoSuccessEffect}{0.0838}
\newcommand{\SPTwoSuccessCILow}{0.0763}
\newcommand{\SPTwoSuccessCIHigh}{0.0915}
\newcommand{\SPTwoIncompleteEffect}{-0.0982}
\newcommand{\SPTwoAlignmentEffect}{0.0877}
\newcommand{\SPTwoPDEffect}{-0.119}
\newcommand{\SPTwoPDCILow}{-0.127}
\newcommand{\SPTwoPDCIHigh}{-0.111}
\newcommand{\SPTwoPDPHolm}{1{,}000}


\subsection{SP2: servicio operacional heterogéneo}
\label{subsec:sp2}

SP2 conserva exclusividad, cierre y poses estáticas, pero cada par aporta distinto servicio. El índice combina carga útil nominal, batería y distancia sin representar mecánica. Se caracteriza la alineación marginal y se comparan puntuaciones. SP3 introduce factibilidad física.

\subsubsection{Índice operacional de servicio}

La Figura~\ref{fig:sp2-scenario} compara coaliciones de igual cardinalidad. Solo una supera diez unidades por carga nominal, batería y distancia.

\begin{figure}[!htbp]
  \centering
  \begin{spdiagram}
    \spdrawarena{SP2 \textbar\ misma cardinalidad, distinto servicio}
    \node[spload,minimum width=1.55cm,minimum height=0.88cm] (load) at (6.65, 2.75) {$L_k$};
    \node[spannotation] at (6.65, 1.78) {$d_k^{\mathrm{srv}}=10$};
    \node[sprobot] (r1) at (1.35, 4.10) {$R_1$};
    \node[sprobot] (r2) at (2.85, 3.00) {$R_2$};
    \node[spannotation,anchor=west] at (0.55, 3.45) {$c_1^{\mathrm{pay}}=8$, $b_1=0.60$, $\delta_{1k}=14\,\mathrm m$};
    \node[spannotation,anchor=west] at (0.55, 2.30) {$c_2^{\mathrm{pay}}=7$, $b_2=0.55$, $\delta_{2k}=11\,\mathrm m$};
    \draw[spcandidate] (r1) -- (load);
    \draw[spcandidate] (r2) -- (load);
    \node[spannotation,text=SPRed] at (3.10, 1.25) {$C_k^A:\ 4+4=8<10$};
    \node[sprobotgreen] (r3) at (10.45, 4.10) {$R_3$};
    \node[sprobotgreen] (r4) at (11.95, 3.00) {$R_4$};
    \node[spannotation,anchor=east] at (12.75, 3.45) {$c_3^{\mathrm{pay}}=10$, $b_3=0.90$, $\delta_{3k}=6\,\mathrm m$};
    \node[spannotation,anchor=east] at (12.75, 2.30) {$c_4^{\mathrm{pay}}=8$, $b_4=0.85$, $\delta_{4k}=5\,\mathrm m$};
    \draw[spassign] (r3) -- (load);
    \draw[spassign] (r4) -- (load);
    \node[spannotation,text=SPGreen!70!black] at (10.40, 1.25) {$C_k^B:\ 7+5=12\geq10$};
    \draw[spcandidate] (0.60, 0.42) -- (1.45, 0.42);
    \node[splegend,anchor=west] at (1.58, 0.42) {alternativa insuficiente};
    \draw[spassign] (4.48, 0.42) -- (5.33, 0.42);
    \node[splegend,anchor=west] at (5.46, 0.42) {asignación que completa};
    \node[splegend,anchor=west] at (9.00, 0.42) {$e_{ik}$ en unidades de servicio};
  \end{spdiagram}
  \caption{Vista conceptual de SP2: dos coaliciones con cardinalidad dos y distinto servicio operacional.}
  \label{fig:sp2-scenario}
  \viuownsource
\end{figure}

Sea $c_i^{\mathrm{pay}}$ la carga útil nominal del robot $i$ en kilogramos, $b_i\in[0, 1]$ su fracción de batería y $b_i^{\mathrm{res}}\in[0, 1)$ la reserva mínima. La distancia inicial robot--carga es $\delta_{ik}$ y $\ell_d>0$ es la escala operacional de distancia. Se definen los factores adimensionales

\begin{equation}
\label{eq:sp2-discounts}
  \psi_i^b=\max\!\left\{0,\frac{b_i-b_i^{\mathrm{res}}}{1-b_i^{\mathrm{res}}}\right\},
  \qquad
  \psi_{ik}^d=\exp\!\left(-\frac{\delta_{ik}}{\ell_d}\right),
  \qquad
  a_{ik}=\psi_i^b\psi_{ik}^d,
\end{equation}

y, con una escala de referencia $c^{\mathrm{ref}}=1\,\mathrm{kg}$, la contribución de servicio

\begin{equation}
\label{eq:sp2-effective-capacity}
  \begin{aligned}
    e_{ik}&=\frac{c_i^{\mathrm{pay}}}{c^{\mathrm{ref}}}a_{ik},
    & S_k(x)&=\sum_{i=1}^{N}e_{ik}x_{ik},\\
    S_k(x)&\geq d_k^{\mathrm{srv}}
    &&\Longleftrightarrow \text{ umbral operacional cubierto}.
  \end{aligned}
\end{equation}

$a_{ik}$ es disponibilidad y $e_{ik}$ servicio normalizado. $c^{\mathrm{ref}}=1\,\mathrm{kg}$ conserva la escala adimensional sin convertir $e_{ik}$ en masa. La distancia penaliza la llegada y deja intacta la capacidad mecánica. Restricciones de soporte y balance energético serían independientes y no se evaluaron. La compatibilidad se fijó a uno.

\subsubsection{Problemas formales de referencia}

Cobertura parcial y completitud son estimandos distintos. Dos MILP globales usan $x_{ik}\in\{0,1\}$ y servicio cubierto $s_k\in[0,d_k^{\mathrm{srv}}]$, con exclusividad por robot.

La referencia de cobertura normalizada es

\begin{equation}
\label{eq:sp2-capacity-oracle}
\begin{aligned}
  \max_{x,s}\quad
  & \sum_{k=1}^{K}\frac{s_k}{d_k^{\mathrm{srv}}}
    -10^{-5}\sum_{i,k}\delta_{ik}x_{ik} \\
  \text{sujeto a}\quad
  & \sum_k x_{ik}\leq1, && \forall i,\\
  & 0\leq s_k\leq d_k^{\mathrm{srv}},\quad s_k\leq\sum_i e_{ik}x_{ik}, && \forall k,\\
  & x_{ik}\in\{0, 1\}. &&
\end{aligned}
\end{equation}

$10^{-5}$ puntos/m solo desempata. \eqref{eq:sp2-capacity-oracle} maximiza razones de cobertura sin representar masa ni completitud. Su \emph{diferencia de cobertura} no acota otras métricas.

La referencia de puntuación añade $z_k\in\{0, 1\}$, activo solo si se cumplen demanda y cardinalidad $n_k$, y usa los pesos del experimento:

\begin{equation}
\label{eq:sp2-score-oracle}
\begin{aligned}
  \max_{x,z,s}\quad
  & 5\sum_k V_kz_k
    +100\sum_k\frac{s_k}{d_k^{\mathrm{srv}}}
    -5\times10^{-4}\sum_{i,k}\delta_{ik}x_{ik}
    -10^{-5}\sum_{i,k}E_{ik}^{\mathrm{trav}}x_{ik} \\
  \text{sujeto a}\quad
  & \sum_kx_{ik}\leq1, && \forall i,\\
  & \sum_i e_{ik}x_{ik}\geq d_k^{\mathrm{srv}}z_k, && \forall k,\\
  & \sum_i x_{ik}\geq n_kz_k, && \forall k,\\
  & 0\leq s_k\leq d_k^{\mathrm{srv}},\quad s_k\leq\sum_i e_{ik}x_{ik}, && \forall k,\\
  & x_{ik},z_k\in\{0, 1\}. &&
\end{aligned}
\end{equation}

$V_k$ es adimensional y $E_{ik}^{\mathrm{trav}}$ se mide en Wh, con coeficientes inversos. Ambas referencias tienen ventaja global. Solo un certificado permite llamarlas exactas.

\subsubsection{Métodos de comparación}

Los MILP separan cobertura y completitud. Se comparan reglas voraces, poblacionales, primal--dual, lineales y neuronales. Húngaro expandido y CBBA-capacidad son adaptaciones sin garantías heredadas \parencite{kuhn1955Hungarian,choiBrunetHow2009CBBA}. En la Tabla~\ref{tab:sp2-methods}, $S$ son puestos, $d$ rasgos y $P$ pesos.

\begin{table}[!htbp]
  \centering
  \caption{Comparación computacional de métodos en SP2.}
  \label{tab:sp2-methods}
  \viutablefont
  \setlength{\tabcolsep}{1.45pt}
  \renewcommand{\arraystretch}{1.02}
  \begin{tabularx}{\textwidth}{>{\raggedright\arraybackslash}p{2.1cm} >{\raggedright\arraybackslash}p{1.55cm} >{\raggedright\arraybackslash}p{2.25cm} >{\raggedright\arraybackslash}p{2.0cm} >{\raggedright\arraybackslash}p{1.8cm} Y}
    \toprule
    Método/fuente & Clase & Coste & Arq./información & Parámetros & Garantía/límite \\
    \midrule
    MILP de puntuación & 0--1 MILP; NP-difícil por la selección opcional heredada de SP1 & Exponencial en peor caso; optimizador con límite temporal y brecha & Central; todos los pares, demandas, costes y recompensas & Tres pesos; límite temporal & Óptimo solo con certificado; objetivo distinto de cobertura. \\
    MILP de cobertura & 0--1 MILP; clase específica no demostrada & Exponencial en el peor caso del optimizador & Central; capacidades efectivas y demandas globales & Dos pesos; límite temporal & Techo de cobertura. No maximiza completitud. \\
    Húngaro expandido \parencite{kuhn1955Hungarian} & Aproximación LAP en P & $O(M^3)$, $M=\max\{N,S\}$ & Central; matriz robot--puesto global & Número y regla de puestos & Exacto solo para la aproximación. No es una reducción exacta de capacidad. \\
    Voraz de capacidad & Problema entero NP-difícil; heurística & $O(NK\log(NK))$ si se ordena una vez & Local/secuencial con candidatos disponibles & Desempate & Sin cota; sensible al orden y al déficit visible. \\
    Aproximación CBBA-capacidad \parencite{choiBrunetHow2009CBBA} & MRTA con capacidades; heurística & $O(NK)$ por barrido de pujas más consenso & Pujas y conflictos entre vecinos & Tres parámetros de puntuación & La puja modificada impide transferir la garantía CBBA. \\
    Replicator/BNN/Smith \parencite{sandholm2010population} & Ordenación secuencial, sin EDO & $O(NK)$ por matriz de puntuaciones & Déficit agregado o visibilidad limitada & Cuatro/cinco parámetros ajustados & No transfiere invariancia ni convergencia de las dinámicas continuas. \\
    Primal--dual global/local & Relajación continua; cierre entero heurístico & $O(NK)$ por actualización de puntuación/precio & Global agregado o vecinal local & Cinco ganancias/pasos & No aporta certificado primal--dual del cierre entero. \\
    Imitación lineal / red neuronal & Inferencia polinómica & $O(NKd)$ / $O(NKP)$; $d=8$, $P=\SPTwoNeuralParameters$ & Rasgos robot--carga durante decodificación & 8 o $P$ pesos; 1/3 hiperparámetros & Sin garantía fuera de la distribución ensayada ni optimalidad. \\
    \bottomrule
  \end{tabularx}
  \viusource{elaboración propia a partir de los métodos ejecutados y las fuentes citadas}
\end{table}

\subsubsection{Juego de capacidad y alineación marginal}

\begin{contribucion}{puntuación de servicio y potencial marginal de SP2}
Para aislar la propiedad estratégica, sea $\rho_{ik}\in[0, 1]$ la preferencia continua del robot $i$ por la carga $k$, con $\sum_k\rho_{ik}\leq1$, y manténgase fija la matriz $E=[e_{ik}]$ durante el instante de decisión. Se usa la presión de déficit lineal saturada $\sigma(r)=[1-r]_+$ sobre la razón $r_k=S_k(\rho)/d_k^{\mathrm{srv}}$. Si $g_{ik}$ es un coste normalizado, las puntuaciones plana y marginal son

\begin{equation}
\label{eq:sp2-payoffs}
  f_{ik}^{\mathrm{plain}}
  =V_k\left[1-\frac{S_k(\rho)}{d_k^{\mathrm{srv}}}\right]_+-g_{ik},
  \qquad
  f_{ik}^{\mathrm{marg}}
  =\frac{e_{ik}}{d_k^{\mathrm{srv}}}V_k\left[1-\frac{S_k(\rho)}{d_k^{\mathrm{srv}}}\right]_+-g_{ik}.
\end{equation}

La puntuación plana ignora contribuciones distintas. La marginal pondera la presión por la fracción de demanda cubierta.

\begin{proposicion}[Alineación marginal del servicio operacional]
\label{prop:sp2-marginal-alignment}
En una región donde $S_k(\rho)<d_k^{\mathrm{srv}}$, con $E$, $d_k^{\mathrm{srv}}$, $V_k$ y $g_{ik}$ fijos, el campo marginal de~\eqref{eq:sp2-payoffs} es el gradiente del potencial
\begin{equation}
\label{eq:sp2-potential}
  \Phi(\rho)=
  \sum_k V_k\int_0^{S_k(\rho)/d_k^{\mathrm{srv}}}[1-s]_+\,\mathrm ds
  -\sum_{i,k}g_{ik}\rho_{ik}.
\end{equation}
El campo plano no es, en general, integrable como potencial exacto cuando existen $i,j,k$ con $e_{ik}\neq e_{jk}$. La conclusión no implica cierre entero, optimalidad combinatoria ni convergencia de una discretización concreta.
\end{proposicion}
\end{contribucion}

La regla de la cadena prueba la alineación. La puntuación plana pierde simetría cruzada con heterogeneidad. El Anexo~\ref{app:sp2-marginal-proof} incluye la frontera saturada.

Un robot requeriría $e_{ik}$, $(d_k^{\mathrm{srv}},V_k)$ y una estimación de $S_k$. Casi todas las reglas usaron déficit global. Solo primal--dual limitó radio. El experimento aísla la puntuación y deja pendiente el estimador. Las variantes Smith fueron ordenaciones, no EDO.

\subsubsection{Regulación del déficit de capacidad}

La salida que SP2 intenta regular ya no es un conteo, sino el déficit relativo de servicio:
\begin{equation}
 z_{2,k}(\rho)=
 \left[1-\frac{S_k(\rho)}{d_k^{\mathrm{srv}}}\right]_+,
 \qquad z_{2,k}^{\mathrm{ref}}=0
 \label{eq:sp2-strategic-regulation}
\end{equation}
para cada carga seleccionada. Si no pueden cubrirse todas, el optimizador selecciona cuáles regulan $z_{2,k}=0$ y las demás se liberan.

La realimentación mínima del robot $i$ está formada por su contribución $e_{ik}$, su coste $g_{ik}$, los parámetros públicos $(d_k^{\mathrm{srv}},V_k)$ y una estimación $\widehat S_{ik}$ del servicio agregado. Con esos datos calcula
\begin{equation}
 \widehat f_{ik}^{\mathrm{marg}}
 =\frac{e_{ik}}{d_k^{\mathrm{srv}}}V_k
  \left[1-\frac{\widehat S_{ik}}{d_k^{\mathrm{srv}}}\right]_+-g_{ik},
 \qquad
 \rho_i^{m+1}=P_{\Delta_i}\!\left(
 \rho_i^m+\Delta t\,F_i(\widehat f_i^m)\right),
 \label{eq:sp2-feedback-candidate}
\end{equation}
donde $F_i$ puede ser Smith, BNN, replicator o gradiente proyectado. La Proposición~\ref{prop:sp2-marginal-alignment} da estructura de gradiente con agregados exactos y $E$ fija, no estabilidad, convergencia discreta ni preservación tras cierre.

El experimento no ejecutó~\eqref{eq:sp2-feedback-candidate}: ordenó secuencialmente y fijó $\widehat S_{ik}=S_k$ casi siempre. Posiciones, batería y $E$ son estáticos. El experimento no incluye una planta. SP3 introduce \emph{wrench} y SP4 transporte.
\subsubsection{Diseño experimental}

La comparación utiliza las semillas $3100$--$3139$ y cinco familias de escenarios: ligero mixto, capacidad balanceada, capacidad pesada, batería limitada y Monte Carlo. Las variaciones internas producen $\SPTwoWorlds$ mundos pareados por método. Se evaluaron once métodos y dos referencias, para $\SPTwoComparisonRows$ observaciones método--mundo. Todas cumplieron las comprobaciones de capacidad, asignación, finitud y consistencia. La ablación plano--marginal añade cuatro reglas y las dos referencias sobre los mismos mundos, con $\SPTwoAblationRows$ comprobaciones sin fallo.

Se miden cobertura, completitud, brechas frente a ambos MILP, servicio incompleto, alineación, distancia, energía, mensajes y CPU. Por mundo se usan Friedman/Kendall, Wilcoxon, media pareada con \mbox{bootstrap} y Holm. Se preespecificaron los efectos de la puntuación marginal y primal--dual frente a voraz.

Las observaciones permiten repetir análisis y figuras, pero no generar mundos nuevos. Esta carencia limita la simulación y conserva la reproducibilidad del recálculo presentado.

\subsubsection{Resultados}

La referencia de cobertura alcanzó la mayor cobertura media, $\SPTwoCoverageCoverage$, pero completó $\SPTwoCoverageSuccess$ de las cargas. El oráculo de puntuación obtuvo una cobertura algo menor, $\SPTwoScoreCoverage$, y una completitud de $\SPTwoScoreSuccess$. La inversión confirma que repartir servicio parcial y completar tareas son objetivos diferentes (Tabla~\ref{tab:sp2-main-results}).

\begin{table}[!htbp]
  \centering
  \caption[SP2 en 1560 mundos pareados. Diferencia de cobertura frente al oráculo de cobertura.]{SP2 en $\SPTwoWorlds$ mundos pareados. ``Dif. cobertura'' usa~\eqref{eq:sp2-capacity-oracle}, no una brecha universal.}
  \label{tab:sp2-main-results}
  \resizebox{\textwidth}{!}{\begin{tabular}{lrrrrr}
\toprule
Método & Cobertura & Completitud & Brecha & Dif. cobertura & CPU [ms] \\
\midrule
Oráculo de puntuación & 0.514 & 0.487 & 0.000 & 0.029 & 40.84 \\
Referencia de cobertura & 0.529 & 0.227 & 0.049 & 0.000 & 11.31 \\
Puntuación neuronal & 0.507 & 0.263 & 0.109 & 0.050 & 1.01 \\
Imitación lineal & 0.514 & 0.218 & 0.118 & 0.038 & 0.98 \\
Voraz de servicio & 0.510 & 0.219 & 0.276 & 0.044 & 0.30 \\
Puntuación inspirada en replicator & 0.525 & 0.139 & 0.297 & 0.026 & 0.36 \\
Puntuación inspirada en Smith & 0.521 & 0.103 & 0.343 & 0.033 & 0.35 \\
\bottomrule
\end{tabular}
}
  \viusource{elaboración propia a partir de los datos procesados de SP2}
\end{table}

% Gráfico generado conservado en el artefacto procesado; las tablas siguientes
% contienen los estimandos y contrastes exactos.
\iffalse
\begin{figure}[!htbp]
  \centering
  \includegraphics[width=\textwidth]{../results/processed/sp2/SP2_EFFECTIVE_CAPACITY_EVIDENCE_v1/figures/fig-sp2-effective-capacity.pdf}
  \caption{Resultados de SP2: cobertura frente a completitud, calidad frente a coste y ablación de la puntuación plana frente a la marginal.}
  \label{fig:sp2-results}
  \viuownsource
\end{figure}
\fi

Entre los métodos aproximados, la red neuronal obtuvo la menor brecha de puntuación, $\SPTwoNeuralGap$, con completitud $\SPTwoNeuralSuccess$ y $\SPTwoNeuralRuntime\,\mathrm{ms}$ de CPU media. Smith fue más barata, $\SPTwoSmithRuntime\,\mathrm{ms}$, pero alcanzó una brecha de $\SPTwoSmithGap$ y completitud $\SPTwoSmithSuccess$. La regla interpretable reduce el coste computacional, aunque no encabeza la calidad de asignación.

La ablación de la Tabla~\ref{tab:sp2-ablation} concuerda con la Proposición~\ref{prop:sp2-marginal-alignment} dentro del asignador secuencial. La corrección marginal elevó la completitud de $\SPTwoPlainSuccess$ a $\SPTwoMarginalSuccess$ y redujo la brecha de $\SPTwoPlainGap$ a $\SPTwoMarginalGap$. El efecto pareado sobre la brecha fue $\SPTwoGapEffect$ (IC del 95\,\% $[\SPTwoGapCILow,\SPTwoGapCIHigh]$, $p_{\mathrm{Holm}}=\SPTwoGapPHolm$). La completitud aumentó $\SPTwoSuccessEffect$ (IC $[\SPTwoSuccessCILow,\SPTwoSuccessCIHigh]$). El servicio atrapado en cargas incompletas disminuyó $\SPTwoIncompleteEffect$. Estos cambios respaldan la ponderación $e_{ik}/d_k^{\mathrm{srv}}$, no una trayectoria Smith que no se ejecutó.

\begin{table}[!htbp]
  \centering
  \caption{Ablación pareada de las puntuaciones plana y marginal en SP2.}
  \label{tab:sp2-ablation}
  \resizebox{0.82\textwidth}{!}{\begin{tabular}{lrrrr}
\toprule
Variante & Completitud & Brecha & Serv. incompleto & Alineación \\
\midrule
Puntuación plana & 0.051 & 0.372 & 0.946 & 0.049 \\
Puntuación marginal & 0.135 & 0.157 & 0.848 & 0.137 \\
\bottomrule
\end{tabular}
}
  \viusource{elaboración propia a partir de los datos procesados de SP2}
\end{table}

Primal--dual no superó al método voraz: la diferencia media de completitud fue $\SPTwoPDEffect$, IC del 95\,\% $[\SPTwoPDCILow,\SPTwoPDCIHigh]$, con $p_{\mathrm{Holm}}=\SPTwoPDPHolm$. Representar el déficit mediante un precio no basta si la selección dispersa capacidad o si la puntuación local no respeta la prioridad todo-o-nada. La ponderación marginal mejora la regla ensayada. El resultado no se extiende a cualquier método primal--dual.

El índice de servicio combina disponibilidad y llegada, pero no certifica capacidad física, energía de misión, fuerza ni torque. La mayoría de los métodos consulta el déficit agregado y no prueba un estimador vecinal. Las variantes poblacionales son ordenaciones secuenciales, no integraciones de sus EDO. El alcance se limita a los mundos evaluados.

\subsubsection{Síntesis}

SP2 establece que una cardinalidad correcta no garantiza cobertura operacional y que una cobertura alta tampoco garantiza completitud. La Proposición~\ref{prop:sp2-marginal-alignment} demuestra, para $E$ fija, que ponderar por $e_{ik}/d_k^{\mathrm{srv}}$ restaura la estructura de gradiente perdida por la puntuación plana. La ablación concentra mejor el servicio con esa corrección. La red neuronal conserva la menor brecha entre los métodos no oráculo y primal--dual no supera al método voraz.

$S_k\geq d_k^{\mathrm{srv}}$ no certifica soporte, contacto, dirección de fuerza ni torque. SP3 sustituye este índice por puestos, geometría y \emph{wrench} admisible.
